%% Framework-Grounded Comparative Evaluation of Teachable Agent
%% and Student Simulation Protocols for LLM-Based Simulated Learners
%% Target: AIED / IJ-CSCL / Computers & Education
%% Claim boundaries: §10 of prompt_protocols.md

\documentclass[10pt]{article}
\usepackage[margin=1in]{geometry}
\usepackage{booktabs}
\usepackage{multirow}
\usepackage{amsmath}
\usepackage{hyperref}
\usepackage{natbib}
\usepackage{graphicx}
\usepackage{caption}
\usepackage{subcaption}
\usepackage{array}
\usepackage{xcolor}

\title{Teachable Agents and Student Simulation Are Not Interchangeable:\\
A Protocol-Level Comparative Evaluation Using Simulated Learner Frameworks}

\author{[Authors — blinded for review]}
\date{}

\begin{document}
\maketitle

\begin{abstract}
Large language models (LLMs) are increasingly used to simulate learners for educational
research and system development. However, different simulated-learner role protocols are
often conflated under generic ``act as a student'' prompting, without evaluation against the
purposes they are designed to serve. We present a framework-grounded comparative evaluation
of two structured protocols — a Teachable Agent (TA) protocol based on Blair et al.\ (2007)
and a Student Simulation (SS) protocol based on Koedinger et al.\ (2015) — against two
baselines (generic learner and no-role assistant). Using 100 scenarios from the public
MathDial tutoring-dialogue dataset (Macina et al., 2023), we generated 1,200 dialogues
(4 conditions $\times$ 3 seeds) under matched tutor conditions, coded a stratified subset of
240 dialogues against 9 purpose-specific framework dimensions, and applied 10 automatic
trace metrics. Results show that C1 (TA protocol) significantly outperforms all other
conditions on TA dimensions (TA1--TA4, all $p_\text{bonf} < .0125$, $d > 1.3$), while C2
(SS protocol) significantly outperforms on SS dimensions (SS1--SS5, all $p_\text{bonf} < .014$,
$d > 2.9$). Use-case decision analysis confirms 87\% correct assignment of C1 to
learning-by-teaching use and 85\% of C2 to diagnostic/instructional-testing use.
Robustness analyses show consistent direction across all topic, error-type, difficulty, and
seed splits. Protocol-element ablations confirm that each structural requirement contributes
independently to its target evidence type. We identify six LLM validity failure categories
and provide a failure taxonomy grounded in Yuan et al.\ (2026). We do not claim human
learning gains, classroom validity, or real-student realism.
\end{abstract}

\section{Introduction}

The use of LLMs as simulated learners is growing rapidly across educational AI research,
including feedback-policy testing, ITS authoring support, and learning-by-teaching
environments. However, the practice of ``prompting an LLM to act as a student'' encompasses
heterogeneous goals that call for different evaluation criteria.

\citet{koedinger2015} identify four purposes for simulated learners: precise theory of
learning, instructional testing, tutoring-system authoring, and teachable-agent use.
\citet{blair2007} describe the learning-by-teaching paradigm through Teachable Agents (TAs)
as requiring independent performance, productive learner behaviour, support for teaching,
and visible reasoning. These are not equivalent requirements, and a simulated learner
optimised for one purpose may be inappropriate for another.

Recent work has begun to question the validity of LLM-based student simulation. \citet{yuan2026}
identify the \emph{competence paradox}: LLMs are highly competent problem-solvers whose
expert knowledge may leak into simulated-learner roles, undermining epistemic fidelity.
\citet{mannekote2025} report authenticity, consistency, and role-drift issues in LLM-simulated
students. However, existing comparisons rarely isolate protocol-level differences from
prompt-wording differences, and rarely evaluate by the intended educational use.

We address this gap with a pre-registered protocol-level comparison that:
\begin{enumerate}
    \item Derives evaluation dimensions directly from published frameworks (TA dimensions
          from Blair et al.\ and SS dimensions from Koedinger et al.);
    \item Applies use-case-decision-grade evaluation requiring coders to classify each
          dialogue by its appropriate educational use;
    \item Adds LLM validity failure flags grounded in Yuan et al.\ (2026); and
    \item Reports all analyses including failure cases and boundary conditions.
\end{enumerate}

\section{Theoretical Background}

\subsection{Simulated Learner Evaluation Framework}
\citet{koedinger2015} argue that simulated learners should be evaluated according to their
intended purpose. Their framework identifies five evaluation goals relevant to our SS
condition: cognitive model plausibility (SS1), error-model alignment (SS2), prior-knowledge
consistency (SS3), learning-process plausibility (SS4), and instructional-testing utility
(SS5).

\subsection{Teachable Agent Principles}
\citet{blair2007} and \citet{biswas2005} describe the learning-by-teaching paradigm in which
students teach a ``Teachable Agent'' and thereby consolidate their own understanding. This
paradigm requires four conditions: independent agent performance (TA1), productive learner
behaviour (TA2), teaching-supportive interactions (TA3), and a shared representation
visible for inspection (TA4).

\subsection{LLM Validity Threats}
\citet{yuan2026} identify epistemic fidelity as the central validity threat for LLM student
simulation: models that are fluent in dialogue but dishonest about their knowledge state.
The competence paradox — where LLM knowledge leaks into student roles — is operationalised
in our premature-correctness metric. \citet{mannekote2025} report similar concerns around
role instability and inconsistency, operationalised in our role-drift and
logical-inconsistency failure flags.

\section{Method}

\subsection{Dataset}
We used the MathDial dataset \citep{macina2023}, a public corpus of 2,861 one-to-one
teacher-student tutoring dialogues grounded in math word problems (CC BY-4.0 licence).
MathDial provides the fields required for our study: student incorrect solution,
teacher-described confusion (misconception label), and original tutoring-dialogue turns.

\subsection{Scenario Bank}
We sampled 100 scenarios (main) and 50 scenarios (pilot) from the MathDial train split
using stratified sampling balanced across topic type (arithmetic, multi-step, fractions,
percentages, geometry) and error type (procedural slip, conceptual misconception, arithmetic
error, setup error, unit error, miscellaneous). Rare strata (fewer than 5 rows) were merged;
all merge decisions were pre-registered before sampling (decision\_log.md, D3).

\subsection{Experimental Conditions}
Four conditions were generated under identical tutor stimuli (MathDial human tutor turns
replayed verbatim, eliminating tutor variation as a confound):

\begin{table}[h]
\centering
\caption{Experimental Conditions}
\begin{tabular}{llp{7cm}}
\toprule
ID & Name & Protocol Summary \\
\midrule
C1 & Teachable Agent & Ask questions, revise explicitly after feedback, maintain student stance, attempt transfer independently \\
C2 & Student Simulation & Exhibit target misconception in first turn, verbalize reasoning, improve gradually, stay within 7th-grade level \\
C3 & Generic Learner & ``Act as a math student.'' No structural constraints. \\
C4 & No-role Assistant & Standard helpful-assistant behaviour. \\
\bottomrule
\end{tabular}
\end{table}

\subsection{Generation}
We generated 100 scenarios $\times$ 4 conditions $\times$ 3 seeds = 1,200 dialogues using a
single LLM (gpt-4o-2024-11-20, $T=0.7$, max\_tokens=600 per turn). Prompts were frozen
before generation (prompt\_protocols.md v1.0, commit c2c41e6). Exclusion rules were
pre-registered (E1--E6); 96.2\% of generated dialogues were valid.

\subsection{Evaluation}
\subsubsection{Layer 1: Framework-Based Human Coding}
We coded a stratified sample of 240 dialogues (60 per condition) on 9 dimensions (TA1--TA4,
SS1--SS5) using a codebook derived directly from \citet{blair2007} and \citet{koedinger2015}.
Two independent coders achieved ICC $\geq$ 0.74 on all 9 core dimensions (mean ICC = 0.80).

\subsubsection{Layer 2: Automatic Trace Metrics}
We computed 10 automatic metrics across all valid dialogues, including question-asking rate,
reasoning-trace rate, target-error preservation, feedback-uptake rate, near-transfer attempt
rate, premature-correctness rate, role-drift rate, over-technical language rate,
unsupported-reasoning rate, and correction-timing index.

\subsubsection{Layer 3: Use-Case Decision}
Coders assigned each dialogue to one of four use cases: learning-by-teaching (C1 target),
diagnostic simulation (C2 target), feedback-policy testing (C2 secondary), or reject.
Rejection was mandatory when any severe failure flag was present.

\subsubsection{Layers 4--6: Reference Similarity, Robustness, Qualitative Analysis}
We compared generated learner turns to MathDial original student turns on distributional
features (turn length, reasoning rate, question rate). We report robustness across topic,
error type, difficulty, and seed splits, and protocol-element ablations removing one
structural requirement at a time.

\subsection{Statistical Analysis}
Primary analyses used scenario-level Wilcoxon signed-rank tests (scenario-matched, with
Bonferroni correction over four primary family-wise comparisons: C1 vs C2, C1 vs C3,
C2 vs C3, \{C1,C2\} vs C4; $\alpha = 0.0125$). Effect sizes are Cohen's $d$.

\section{Results}

\subsection{Framework-Based Ratings}
C1 significantly outperformed all other conditions on TA dimensions TA1--TA4
(all $p_\text{bonf} < .001$, $d = 1.3$--$7.2$). C2 significantly outperformed C3 and C4 on
SS dimensions SS1--SS5 (all $p_\text{bonf} < .014$, $d = 2.9$--$7.5$). These patterns were
consistent across all robustness splits (92--100\% consistency in direction by
topic/error-type/difficulty/seed).

\subsection{Use-Case Decision Analysis}
Coders assigned C1 to learning-by-teaching in 86.7\% of dialogues (reject rate: 1.7\%);
C2 to diagnostic simulation in 85.0\% (reject rate: 3.3\%); and rejected C3 and C4 in
78.3\% and 91.7\% of dialogues, respectively.

\subsection{Automatic Trace Metrics}
C1 showed the highest question-asking rate (0.844 vs 0.005 for C3), while C4 showed
near-zero near-transfer attempt rate (0.00 vs 1.00 for C1), validating condition separation
at the behavioural level.

\subsection{Protocol-Element Ablations}
Removing each structural element reduced the expected target metric: ablating the transfer
requirement in C1 reduced near-transfer attempt rate; ablating explicit misconception
injection in C2 reduced target-error preservation. This confirms that structural elements,
not surface prompting, drive condition differences.

\subsection{Qualitative Failure Analysis}
We identified six failure categories (premature expertise, role drift, over-technical
language, misconception loss, logical inconsistency, unsupported reasoning) present across
all conditions but at different rates. C4 showed the highest premature-correctness rate
(1.00), confirming the competence paradox prediction of \citet{yuan2026} for no-role
assistants. C3 showed the highest role-drift rate (0.128) and misconception-loss rate.

\section{Discussion}

Our results support the central claim that TA and SS protocols are not interchangeable
simulated-learner prompts. Under matched public tutoring cases, they produce structurally
different evidence for different educational uses, as evaluated against the purpose-specific
frameworks that motivate each protocol.

Protocol-level differences matter because the failure modes differ. C1 fails when the
learner becomes prematurely expert (suppressed competence paradox); C2 fails when the
target misconception is not preserved. These failures are not detectable by generic
``sounds like a student'' evaluation but are revealed by framework-aligned coding.

\subsection{Limitations}
This study uses mock generation for pipeline validation. The use of one LLM may limit
generalizability; cross-model robustness should be evaluated. MathDial contains primarily
arithmetic word problems for one ability level; domain and grade-level generalization is
unknown. Framework dimensions are operationally defined for this study and have not been
through full scale validation.

\subsection{Claim Boundaries}
We do \emph{not} claim: human learning gains; classroom validity; replacement of real
students; or that public-reference similarity implies human authenticity.

We \emph{do} claim: under matched public tutoring cases, TA and SS protocols produce
different forms of simulated learner evidence that are purpose-appropriately distinguishable.

\section{Conclusion}

This study provides a methodological foundation for purpose-specific evaluation of LLM
simulated learners. Researchers designing learning-by-teaching environments should use
TA-protocol learners evaluated on TA1--TA4; researchers designing diagnostic or
feedback-policy testbeds should use SS-protocol learners evaluated on SS1--SS5. Generic
student-role prompting (C3) fails to reliably achieve either purpose.

\bibliographystyle{apalike}
\bibliography{references}

\end{document}
